% from ega1-10-fr.tex
\oldpage[I]{186}
Si $U$ est une partie ouverte de $\mathfrak{X}$, l'injection canonique
dans $\mathfrak{X}$ du préschéma formel induit par $\mathfrak{X}$ sur $U$
est un morphisme de préschémas formels (et même un
\emph{monomorphisme} d'espaces topologiquement annelés
(\hyperref[0.4.1.1-fr]{0, 4.1.1})).

\begin{proposition}[10.4.6]
\label{I.10.4.6-fr}
Soient $\mathfrak{X}$ un préschéma formel,
$\mathfrak{S}=\operatorname{Spf}(A)$ un schéma formel affine. Il existe
une correspondance biunivoque canonique entre les morphismes du préschéma
formel $\mathfrak{X}$ dans le préschéma formel $\mathfrak{S}$ et les
homomorphismes continus de l'anneau $A$ dans l'anneau topologique
$\Gamma(\mathfrak{X},\mathscr{O}_{\mathfrak{X}})$.
\end{proposition}

La démonstration est la même que celle de
(\hyperref[I.2.2.4-fr]{2.2.4}), en remplaçant «~homomorphisme~» par
«~homomorphisme continu~», «~ouvert affine~» par «~ouvert formel
affine~», et en utilisant (\hyperref[I.10.2.2-fr]{10.2.2}) au lieu de
(\hyperref[I.1.7.3-fr]{1.7.3})~; nous en laissons les détails au lecteur.

\begin{env}[10.4.7]
\label{I.10.4.7-fr}
Étant donné un préschéma formel $\mathfrak{S}$, on dit que la donnée d'un
préschéma formel $\mathfrak{X}$ et d'un morphisme
$\varphi:\mathfrak{X}\to\mathfrak{S}$ définit un préschéma formel
$\mathfrak{X}$ \emph{au-dessus de $\mathfrak{S}$} ou un
\emph{$\mathfrak{S}$-préschéma formel}, $\varphi$ étant appelé le
\emph{morphisme structural} du $\mathfrak{S}$-préschéma
$\mathfrak{X}$. Si $\mathfrak{S}=\operatorname{Spf}(A)$, où $A$ est un
anneau admissible, on dit aussi que le $\mathfrak{S}$-préschéma formel
$\mathfrak{X}$ est un \emph{$A$-préschéma formel} ou un préschéma formel
\emph{au-dessus de $A$}. Un préschéma formel quelconque peut toujours
être considéré comme un préschéma formel au-dessus de $\mathbf{Z}$ (muni
de la topologie discrète).

Si $\mathfrak{X}$, $\mathfrak{Y}$ sont deux $\mathfrak{S}$-préschémas
formels, on dit qu'un morphisme
$u:\mathfrak{X}\to\mathfrak{Y}$ est un
\emph{$\mathfrak{S}$-morphisme} si le diagramme
\[
  \xymatrix{
    \mathfrak{X}\ar[rr]^u\ar[rd] & &
    \mathfrak{Y}\ar[ld]\\
    & \mathfrak{S}
  }
\]
(où les flèches obliques sont les morphismes structuraux) est
commutatif. Avec cette définition, les $\mathfrak{S}$-préschémas formels
(pour $\mathfrak{S}$ fixé) forment une \emph{catégorie}. On désigne par
$\operatorname{Hom}_{\mathfrak{S}}(\mathfrak{X},\mathfrak{Y})$
l'ensemble des $\mathfrak{S}$-morphismes du $\mathfrak{S}$-préschéma
formel $\mathfrak{X}$ dans le $\mathfrak{S}$-préschéma formel
$\mathfrak{Y}$. Lorsque $\mathfrak{S}=\operatorname{Spf}(A)$, on dit
aussi \emph{$A$-morphisme} au lieu de
\emph{$\mathfrak{S}$-morphisme}.
\end{env}

\begin{env}[10.4.8]
\label{I.10.4.8-fr}
Comme tout schéma affine peut être considéré comme un schéma formel affine
(\hyperref[I.10.1.2-fr]{10.1.2}), tout préschéma (usuel) peut être
considéré comme un préschéma formel. Il résulte en outre de
(\hyperref[I.10.4.5-fr]{10.4.5}) que pour les préschémas \emph{usuels},
les morphismes (resp. $S$-morphismes) de préschémas \emph{formels}
coïncident avec les morphismes (resp. $S$-morphismes) définis au
\S~\hyperref[section:I.2-fr]{2}.
\end{env}

\subsection{Idéaux de définition des préschémas formels.}
\label{subsection:I.10.5-fr}

\begin{env}[10.5.1]
\label{I.10.5.1-fr}
Soit $\mathfrak{X}$ un préschéma formel~; on dit qu'un
$\mathscr{O}_{\mathfrak{X}}$-Idéal $\mathscr{J}$ est un
\emph{faisceau d'idéaux de définition} (ou un
\emph{Idéal de définition}) de $\mathfrak{X}$ si tout
$x\in\mathfrak{X}$ possède un voisinage ouvert formel affine $U$ tel que
$\mathscr{J}|U$ soit un Idéal de définition du schéma formel affine
induit par $\mathfrak{X}$ sur $U$
(\hyperref[I.10.3.3-fr]{10.3.3})~; en vertu de
(\hyperref[I.10.3.1.1-fr]{10.3.1.1}) et
(\hyperref[I.10.4.3-fr]{10.4.3}), pour tout ouvert
$V\subset\mathfrak{X}$, $\mathscr{J}|V$ est alors un Idéal de définition
du préschéma formel induit par $\mathfrak{X}$ sur $V$.

On dit qu'une famille $(\mathscr{J}_\lambda)$ d'Idéaux de définition de
$\mathfrak{X}$ est un \emph{système fondamental}
